%% Acknowledgements
\chapter*{Acknowledgements}
\addcontentsline{toc}{chapter}{Acknowledgements}

The author gratefully acknowledges the following organisations and individuals whose data, expertise, and frameworks underpin this analysis.

\textbf{Natural England} for the publication and maintenance of the Priority Habitat Inventory (v2.3), the Biodiversity Metric~4.0 calculation tool and supporting technical guidance, and the Magic Map environmental data service.

\textbf{Department for Environment, Food \& Rural Affairs (DEFRA)} for the Statutory Biodiversity Metric~4.0 User Guide and Technical Supplement, the Biodiversity Gain Site Register, and policy guidance on biodiversity net gain under the Environment Act 2021.

\textbf{Ordnance Survey} for access to OS MasterMap Topography Layer, OS Terrain~50, and OS Open Data products via the OS Data Hub API, used under the Public Sector Geospatial Agreement.

\textbf{European Space Agency (ESA)} for the Copernicus Sentinel-2 multispectral imagery archive, accessed via Google Earth Engine, which forms the primary remote sensing data source for both the 2016 and 2026 classification epochs.

\textbf{United States Geological Survey (USGS)} for Landsat~8 OLI/TIRS data products, used as complementary spectral evidence during classification training and validation.

\textbf{Google Earth Engine (GEE)} team for providing the cloud computing platform used for imagery compositing, spectral index calculation, and initial classification prototyping.

\textbf{Derbyshire County Council} for the emerging Local Nature Recovery Strategy (LNRS) data layers and planning constraint information pertaining to the Ashleyhay study area.

\textbf{Derbyshire Wildlife Trust} for publicly available habitat survey guidance relevant to the Peak Fringe landscape character area.

\textbf{University of Derby} for the undergraduate research thesis data on habitat condition in the Ashleyhay area, which provided an independent validation dataset for the cross-examination presented in Chapter~\ref{ch:spotcheck}.

The computational pipeline was developed and executed by Dr~John O'Hare using the DreamLab AI agent-swarm infrastructure. All errors and interpretations remain the responsibility of the author.
